\section*{ID6807: From Linear Modes to Nonlinear Filaments: Zel’dovich Approximation: ∂\_i∂\_jΦ}
\paragraph{Metadata.}
PiID: 3643435439478 | RTEID: RTE:P26433.0927:T5.8545 | UUID: 5f8327ba-e4ea-529e-9953-dad6bd817782

\paragraph{Canonical Statement.}
Now, how do the spherical-Bessel modes feed into \textbackslash{}(\textbackslash{}lambda\_i\textbackslash{})? The tidal tensor is second derivatives of the potential, and the potential itself links to δ via Poisson: \textbackslash{}[ \textbackslash{}nabla\textasciicircum{}2\textbackslash{}Phi(\textbackslash{}mathbf\{q\}) \textbackslash{}propto \textbackslash{}delta(\textbackslash{}mathbf\{q\}). \textbackslash{}] Express \textbackslash{}(\textbackslash{}Phi\textbackslash{}) in the same basis (1), then compute \textbackslash{}(\textbackslash{}partial\_i\textbackslash{}partial\_j\textbackslash{}Phi\textbackslash{}). The angular dependence of the spherical harmonics selects preferred directions where eigenvalues align — those directions become preferred axes of collapse. If the angular superposition produces a strong directional gradient pattern (coherent sign of \textbackslash{}(Y\_\{\textbackslash{}ell\}\textasciicircum{}m\textbackslash{}) across modes), then two eigenvalues will often cross threshold near the same location → **filament formation along the remaining axis**.

\paragraph{Canonical Equation.}
\[
\partial_i\partial_j\Phi
\]

\paragraph{Dictionary.}
Relation: ∂\_i∂\_jΦ

\paragraph{Encyclopedia.}
From Linear Modes to Nonlinear Filaments: Zel’dovich Approximation: ∂\_i∂\_jΦ is a differential equation, written ∂\_i∂\_jΦ. It relates the quantity to its derivatives, governing how the system evolves in space and time. It is built from Φ, j. Within the Trawin Zero Point registry it serves as one element of the framework's mathematical narrative.
